\subsection{把“北方”还给不同的地点}

“北方”常被用作宋朝叙事里一个模糊的方向，好像从开封向北就会进入同一种辽地或同一种边患。实际上，潢河流域、燕云城镇、东北山林、河套灌区、河西绿洲与陕西沿边，有不同的水、土、路和制度接口。辽的上京与南京之间要靠人、物和命令保持联系；西夏的兴庆与黑水城、河西诸城之间也有各自的路网与保存条件。将这些地点放回历史，不是为了绘制精确的现代式国界，而是为了让政治选择重新获得距离和成本。

地图上的约略边界只能帮助读者想象关系，不应替代证据。军事控制会沿城、关、河、道路呈现不均匀的强弱，商路与宗教网络又可能跨过同一条边线。辽、宋、西夏的统治者都需要通过地方中介把命令与资源送到远处；远处的人也会用市场、亲属、寺院和迁徙应对这些要求。所谓多政权时代，正是许多不同强度的接口同时存在，而不是几种颜色整齐相接。

\subsection{制度得失：用可见的接口代替抽象的强弱}

辽的多中心安排使草原、东北和燕云能在相当长时间内被纳入同一王朝，却持续要求宫廷、后族、官署、营卫和城市资源相互协调。它的优势在于不必把所有地区强行改成同一形式；代价是当继承、财政或边防同时紧张时，中心之间的联结更容易断裂。西夏依托水利城镇、军政首领、文字和宗教网络建立国家，能够在宋辽金之间保有行动空间；代价则是维持渠系、城寨、文书与军队的资源负担会不断落到有限的地区与人群上。

对两国而言，外交提供的不是摆脱成本的捷径。盟约、册封、贡赐和互市能够为时间紧张的国家争取空间，却要由仓储、道路、翻译和地方社会持续支付。战争会把这些接口暴露为弱点，征服则会迫使新政权重新寻找它们。读者在辽的灭亡和西夏的终局中应保留这一尺度：王朝名号可以突然更换，水渠、寺院、墓地、市场和家族的转换却有不同速度。只有这样，边地才不再是大国年表边缘的空白。
